\section{Bounds: contradiction radius}\label{sec:contradiction}

Theorem~\ref{thm:bounds} returns an interval; a public claim returns a point. The diagnostic measures their distance after standardising each input by its measurement uncertainty.

\begin{definition}[Feasible set]\label{def:feasible}
For inputs $x=(\omega,w,a,f,M,D)$ and exposure range $[\mu_{\mathrm{lo}},\mu_{\mathrm{hi}}]$,
\begin{equation}\label{eq:feasible}
\mathcal F(x)=
\Bigl\{q\in[0,1]:\;
\exists\,\mu\in[\mu_{\mathrm{lo}},\mu_{\mathrm{hi}}]\;\text{with}\;
q=1-\omega\tfrac{w+\mu(a-f)}{w}
\;\text{and}\; qD\le M\Bigr\}.
\end{equation}
\end{definition}

A claim $\hat D_M^{(i)}$ converts to $q^{(i)}=\hat D_M^{(i)}/\hat D$ and is feasible iff $q^{(i)}\in\mathcal F(\hat x)$.

\begin{definition}[Contradiction radius]\label{def:radius}
Let $\hat x=(\hat\omega,\hat w,\hat a,\hat f,\hat M,\hat D)$ collect the measured inputs, related by the accounting identities $\hat a=1-\hat w$ and $\hat f=\hat M/\hat N$, and let $\sigma_\omega,\sigma_w,\sigma_M$ be the uncertainty scales attached to the three independently measured inputs. For a single input $A\in\{\omega,w,M\}$, write $\hat x[A\mapsto A']$ for the input vector with $A$ replaced by $A'$, the identities re-imposed ($a'=1-w'$ when $w$ moves; $f'=M'/\hat N$ when $M$ moves), the revised input confined to its natural domain ($\omega'\in[0,1]$, $w'\in(0,1)$, $M'\ge 0$, and $f'\le a'$), and every other input held at its measured value. The \emph{axis radius} of claim $q^{(i)}$ on axis $A$ is the smallest standardised revision of that input alone that makes the claim feasible,
\begin{equation*}
\rho_A^{(i)}
=\inf\Bigl\{\,\tfrac{|A'-\hat A|}{\sigma_A}
\;:\;q^{(i)}\in \mathcal F\bigl(\hat x[A\mapsto A']\bigr)\Bigr\}
\qquad(\inf\emptyset=+\infty),
\end{equation*}
and the \emph{contradiction radius} of the claim is the cheapest single-input rescue:
\begin{equation}\label{eq:rho}
\rho^{(i)}=\min\bigl\{\rho_\omega^{(i)},\;\rho_w^{(i)},\;\rho_M^{(i)}\bigr\}.
\end{equation}
\end{definition}

\begin{theorem}[Feasibility and contradiction]\label{thm:contradiction}
With $(\hat D,\hat N)$ and the exposure range fixed, a claim is consistent with the independent inputs iff $\rho^{(i)}=0$. If $\rho^{(i)}>0$, then rescuing the claim by revising any single independent input $A$ requires moving it by at least $\rho_A^{(i)}\ge\rho^{(i)}$ standardised units, and the axis attaining the minimum in \eqref{eq:rho} names the cheapest such revision and the institution whose measurement it contradicts.
\end{theorem}

\begin{remark}[Coordinated distortions]\label{rem:joint}
$\rho$ prices the cheapest \emph{single-institution} error. A coordinated distortion of several inputs at once can achieve a smaller largest standardised move by splitting the burden across institutions: replacing the minimum over axes in \eqref{eq:rho} with a joint programme $\inf\max_A |A'-\hat A|/\sigma_A$ over simultaneous revisions gives a value below $\rho$. In the application of Section~\ref{sec:gaza}, granting the disputed claim's upper endpoint simultaneous moves of $\omega$ and $w$ (with $a'=1-w'$ and the census scale $\sigma_w=0.005$ stated there) lowers its radius from $\nRhoOmegaAg{}$ to about $\nRhoJoint{}$: the verdict is unchanged, and the coordinated reading requires two independent institutions' measurements to be wrong at once---exactly the fingerprint the under-fudgible structure is designed to expose.
\end{remark}

\paragraph{Interpretation.} The contradiction radius answers the question a motivated critic must eventually face: \emph{which measurement, exactly, do you say is wrong, and by how much?} A claim with $\rho=0$ is compatible with all independent inputs; nothing more can be said against it on this evidence. A claim with $\rho=2$ requires some input to be off by two standardised units---uncomfortable but conceivable. A claim with $\rho=20$ requires an independently measured quantity to be wrong by twenty times its stated uncertainty scale; since the scales used here are set at or above plausible sampling error, defending such a claim requires asserting a systematic failure or misreporting by a specific institution---and the radius names the institution and the magnitude. (The scale is an input, not an estimated standard deviation; no distributional tail statement is attached to it.) Because $\rho$ standardises each input by its stated uncertainty scale, it is unit-free, and comparable across conflicts and claims to the extent that those scales are comparably defined. It is also symmetric: an implausibly \emph{low} combatant claim generates a positive radius the same way an implausibly high one does. What the radius is not: it is not a posterior probability; it prices the cheapest \emph{single-input} revision, with coordinated multi-input distortions treated in Remark~\ref{rem:joint}; and it inherits any systematic---as opposed to sampling---error in the inputs, which is why Section~\ref{sec:gaza} reports it under every anchor choice.

\begin{figure}[t]
\centering
\begin{tikzpicture}[
  >=Stealth, font=\small,
  vert/.style={circle,draw=black!75,fill=white,inner sep=1.5pt,minimum size=4mm},
  feas/.style={fill=blue!10,draw=blue!55,thick},
  arr/.style={->,>=Stealth,gray!70,thick},
  lab/.style={font=\footnotesize,inner sep=2pt}
]
\coordinate (T)  at ( 0   , 2.5);
\coordinate (BL) at (-3.0 ,-0.9);
\coordinate (BR) at ( 3.0 ,-0.9);

\draw[feas] (T) -- (BL) -- (BR) -- cycle;
\node[lab,align=center] at (0,0.55) {feasible set\\$\mathcal F(\hat x)$};

\node[vert,fill=blue!55] (P) at (0.4, 0.0) {};
\node[lab,below=1pt of P] {feasible $\hat q$};

\node[vert,fill=red!70] (Q) at (4.7,0.95) {};
\node[lab,right=1pt of Q] {disputed $q^{(\mathrm{IDF})}$};

\draw[arr] (Q) -- ($(T)!0.18!(Q)$)  node[midway,above,sloped,font=\scriptsize] {$\rho_\omega$};
\draw[arr] (Q) -- ($(BR)!0.20!(Q)$) node[midway,below,sloped,font=\scriptsize] {$\rho_M$};
\draw[arr] (Q) -- ($(BL)!0.30!(Q)$) node[midway,above,sloped,font=\scriptsize] {$\rho_w$};

\node[lab] at (T)  [above=2pt] {$\hat\omega$ axis (demographic deaths)};
\node[lab] at (BL) [below left=1pt,align=center]  {$\hat w$ axis\\(population)};
\node[lab] at (BR) [below right=1pt,align=center] {$\hat M/D$ axis\\(manpower)};
\end{tikzpicture}
\caption{\textbf{Contradiction triangle (conceptual schematic).} The feasible set $\mathcal F$ is bounded by three independently measured ``axes''; the contradiction radius $\rho=\min(\rho_\omega,\rho_w,\rho_M)$ is the cheapest standardised revision of a single axis that carries a disputed claim into $\mathcal F$ (Definition~\ref{def:radius}). Each axis is owned by a different institution, which is what makes the relation under-fudgible: rescuing a claim on one axis leaves fingerprints on that axis's measurement.}
\label{fig:contradiction-triangle}
\end{figure}

\paragraph{Computation.} Each axis radius is available in closed form. On the demographic axis, inverting \eqref{eq:qmu} gives the composition a claim requires at exposure $\mu$, $\omega_{\mathrm{needed}}(q,\mu)=(1-q)\,\hat w/(\hat w+\mu(\hat a-\hat f))$; granting the claim the most favourable $\mu$ in the admissible range yields the interval of rationalising compositions $[\omega_{\mathrm{needed}}(q,\mu_{\mathrm{hi}}),\,\omega_{\mathrm{needed}}(q,\mu_{\mathrm{lo}})]$, and $\rho_\omega$ is the standardised distance from $\hat\omega$ to that interval (zero if $\hat\omega$ lies inside). The $w$ and $M$ axes are computed the same way through the identities $a'=1-w'$ and $f'=M'/\hat N$. This is the contradiction triangle of Figure~\ref{fig:contradiction-triangle}. In the Gaza application the last two axes cannot rescue the disputed claim at any remotely plausible magnitude (Section~\ref{sec:gaza}, Step~3 quantifies both), so $\rho=\rho_\omega$ there and radii are quoted in units of $\sigma_\omega$.

\paragraph{Diagnostic procedure.}
\begin{enumerate}[label=\textbf{D\arabic*.},leftmargin=*]
\item Fix the analysis window and total deaths $D$; carry the uncertainty in $D$ as a sensitivity dimension, not a constant.
\item Estimate $(w,a)$ from the population baseline and $\omega$ from identified-deaths records, checking each candidate $\omega$ source for selection mechanisms before use.
\item Calibrate the exposure range $[\mu_{\mathrm{lo}},\mu_{\mathrm{hi}}]$ from historical conflicts with known combatant status, and take an independent upper bound on $M$.
\item Compute $\mathcal F(\hat x)$ from \eqref{eq:feasible}.
\item For every disputed $\hat D_M^{(i)}$, report $q^{(i)}=\hat D_M^{(i)}/D$, the membership $q^{(i)}\in\mathcal F$, and $\rho^{(i)}$, under each anchor and each $D$ scenario.
\end{enumerate}
The test is symmetric: it scores ``too many combatants'' and ``zero combatants'' against the same set.
